\section{Conclusion}
\label{Conclusion}

This paper introduces an adaptive CBF safety layer for collaborative autonomous driving that works with the V2Xverse CoDriving stack without retraining perception or planning. By adapting only the safety enforcement while keeping upstream modules fixed, our approach ensures that observed improvements are attributable to the safety layer, suggesting its portability and ease of integration into other collaborative driving platforms.

The approach combines a modular CBF-QP safety filter with a class-aware adaptive $\gamma$ policy trained using constrained PPO and a PID-based Lagrangian. A compact network outputs gains per-step for vehicles, pedestrians, and bicycles from traffic-risk features, while the CBF remains the final mechanism that enforces bounded longitudinal actions. This keeps the safety layer interpretable, lightweight, and directly compatible with an existing collaborative autonomy stack.

When evaluated in Town05 against PID-only and fixed CBF under the same perception, planning, and steering setup, the adaptive method delivered the strongest overall performance: driving score $75.81 \pm 3.35$, collisions per route 0.581, and collisions per kilometer 6.75. Relative to PID-only, collisions per route dropped by 80.3\%, and relative to fixed CBF, they dropped by 25.9\%. The completion of the route remained high at 90.8\%, while the collision-free route rate increased to 59.7\%. The greatest gains on bicycles and pedestrians show that adaptive gain scheduling is especially useful for vulnerable road users in heterogeneous traffic.

Overall, these results support the central claim of the paper: learning to schedule barrier gains online can improve safety without sacrificing route progress. More broadly, they show that meaningful safety improvements can come from a targeted control-layer intervention rather than a full redesign of the autonomy stack.

\textbf{Limitations.} Our formulation does not provide a formal forward-invariance guarantee. The slack variable in the QP admits soft constraint violation when the problem would otherwise be infeasible, and we measure it as active on $30$--$35\%$ of decision steps (Section~\ref{sec:solve-time}); the solver never fell back to the nominal command, but this means the safety property is empirical rather than certified. We bound, rather than eliminate, the risk from the constant-velocity actor assumption via a worst-case acceleration margin (Section~\ref{methods}), and our barrier formulation acts on longitudinal motion only: extending the framework to jointly regulate longitudinal and lateral actions is a natural next step, requiring a coupled vehicle model and richer barrier construction. The checkpoint underlying our headline result predates the checkpoint-selection methodology fix described in Section~\ref{sec:checkpoint-selection}; producing a validation-selected headline result is the top priority for future work. Finally, three additional ablations explicitly called for could not be completed with a trained policy in this revision cycle and remain future work: class-specific-vs-shared-gain and cost-function-weighting variants, and a PID-Lagrangian-vs-plain-Lagrangian comparison -- each requires training a separate policy under the corresponding constraint, which we substituted where possible with training-free proxies on the existing policy (Section~\ref{sec:ablations}) rather than presenting as equivalent evidence. Even with these limitations, the approach is practical as a drop-in safety adapter for existing autonomy stacks and complements continued advances in collaborative perception and planning.
